% Options for packages loaded elsewhere
\PassOptionsToPackage{unicode}{hyperref}
\PassOptionsToPackage{hyphens}{url}
\documentclass[
  11pt,
  a4paper,
]{article}
\usepackage{xcolor}
\usepackage[margin=1in]{geometry}
\usepackage{amsmath,amssymb}
\setcounter{secnumdepth}{-\maxdimen} % remove section numbering
\usepackage{iftex}
\ifPDFTeX
  \usepackage[T1]{fontenc}
  \usepackage[utf8]{inputenc}
  \usepackage{textcomp} % provide euro and other symbols
\else % if luatex or xetex
  \usepackage{unicode-math} % this also loads fontspec
  \defaultfontfeatures{Scale=MatchLowercase}
  \defaultfontfeatures[\rmfamily]{Ligatures=TeX,Scale=1}
\fi
\usepackage{lmodern}
\ifPDFTeX\else
  % xetex/luatex font selection
\fi
% Use upquote if available, for straight quotes in verbatim environments
\IfFileExists{upquote.sty}{\usepackage{upquote}}{}
\IfFileExists{microtype.sty}{% use microtype if available
  \usepackage[]{microtype}
  \UseMicrotypeSet[protrusion]{basicmath} % disable protrusion for tt fonts
}{}
\makeatletter
\@ifundefined{KOMAClassName}{% if non-KOMA class
  \IfFileExists{parskip.sty}{%
    \usepackage{parskip}
  }{% else
    \setlength{\parindent}{0pt}
    \setlength{\parskip}{6pt plus 2pt minus 1pt}}
}{% if KOMA class
  \KOMAoptions{parskip=half}}
\makeatother
\setlength{\emergencystretch}{3em} % prevent overfull lines
\providecommand{\tightlist}{%
  \setlength{\itemsep}{0pt}\setlength{\parskip}{0pt}}
\usepackage{bookmark}
\IfFileExists{xurl.sty}{\usepackage{xurl}}{} % add URL line breaks if available
\urlstyle{same}
\hypersetup{
  pdftitle={Pipeline Improvement Notes},
  hidelinks,
  pdfcreator={LaTeX via pandoc}}

\title{Pipeline Improvement Notes}
\author{}
\date{}

\begin{document}
\maketitle

\section{Pipeline Improvement Notes}\label{pipeline-improvement-notes}

\subsection{Purpose}\label{purpose}

This document is an engineering follow-up note for the current pipeline
revision. It is separate from the submission-facing report and is
intended to capture the most important technical issues observed in the
present build, together with direct, concrete steps to fix them.

The goal is not to redesign the project. The goal is to make targeted
changes that improve the final graph while preserving the strengths of
the current architecture: automated construction, canonical namespace
usage, local-cache reproducibility, provenance capture, and broad source
coverage.

\subsection{Priority 1: Align emitted classes with the generated
ontology}\label{priority-1-align-emitted-classes-with-the-generated-ontology}

\subsubsection{Current issue}\label{current-issue}

The current final graph contains high-frequency auxiliary \texttt{tfl:}
types that are not declared in the generated TBox, including:

\begin{itemize}
\tightlist
\item
  \texttt{tfl:TransportOperator}
\item
  \texttt{tfl:TransitStop}
\item
  \texttt{tfl:Location}
\item
  \texttt{tfl:TransportService}
\item
  \texttt{tfl:TransportEvent}
\end{itemize}

This makes the graph harder to query consistently and weakens the
relationship between the ontology report and the built artifact.

\subsubsection{Direct fix}\label{direct-fix}

\begin{enumerate}
\def\labelenumi{\arabic{enumi}.}
\tightlist
\item
  Open
  \texttt{data\_retrieval/unstructured\_retrieval/triples\_to\_rdf.py}.
\item
  Add a canonical class-alignment map before any \texttt{rdf:type}
  assertions are written.
\item
  Map at least:

  \begin{itemize}
  \tightlist
  \item
    \texttt{TransportOperator\ -\textgreater{}\ TfLOperator}
  \item
    \texttt{TransitStop\ -\textgreater{}\ TfLStation}
  \item
    \texttt{TransportService\ -\textgreater{}\ TfLLine} or a declared
    service class, depending on local usage
  \end{itemize}
\item
  If \texttt{Location} and \texttt{TransportEvent} are still genuinely
  useful, declare them explicitly in \texttt{generateOntology.py};
  otherwise remap or drop them.
\item
  Add a validation script that fails the build if undeclared
  \texttt{tfl:} classes are emitted into the final graph.
\end{enumerate}

\subsubsection{Files to change}\label{files-to-change}

\begin{itemize}
\tightlist
\item
  \texttt{data\_retrieval/unstructured\_retrieval/triples\_to\_rdf.py}
\item
  \texttt{generateOntology.py}
\item
  optionally add \texttt{scripts/check\_ontology\_alignment.py}
\end{itemize}

\subsection{Priority 2: Restore benchmark-oriented instance coverage for
missing competency
questions}\label{priority-2-restore-benchmark-oriented-instance-coverage-for-missing-competency-questions}

\subsubsection{Current issue}\label{current-issue-1}

The graph is broad, but several coursework questions still depend on
tightly targeted instances that are not yet materialized clearly enough
in the final artifact. This is most visible in:

\begin{itemize}
\tightlist
\item
  CQ6 and CQ7 for station facilities
\item
  CQ11 for the peak Oyster fare example
\item
  CQ13 for engineering closures
\item
  CQ15 for the benchmark journey
\item
  CQ20 for fare concessions
\end{itemize}

\subsubsection{Direct fix}\label{direct-fix-1}

\begin{enumerate}
\def\labelenumi{\arabic{enumi}.}
\tightlist
\item
  Create a small post-build enrichment layer containing only
  benchmark-sensitive instances.
\item
  Merge that layer after the automated structured and unstructured
  stages in \texttt{main.py}.
\item
  Keep this layer intentionally small and evidence-backed so it supports
  the queries without turning the build back into a hand-authored graph.
\item
  Name benchmark instances consistently, for example:

  \begin{itemize}
  \tightlist
  \item
    \texttt{tfl:PeakFare\_Adult\_Z1to3}
  \item
    \texttt{tfl:Journey\_BrixtonToCanaryWharf}
  \end{itemize}
\item
  Re-run the full SPARQL harness after each added benchmark layer
  change.
\end{enumerate}

\subsubsection{Files to change}\label{files-to-change-1}

\begin{itemize}
\tightlist
\item
  \texttt{main.py}
\item
  add \texttt{ontologies/manual/benchmark\_enrichment.ttl} or a
  generated equivalent
\item
  optionally add a script that builds benchmark instances from
  structured inputs
\end{itemize}

\subsection{Priority 3: Tighten operator
resolution}\label{priority-3-tighten-operator-resolution}

\subsubsection{Current issue}\label{current-issue-2}

Operator-related queries currently return broad answer sets with
multiple plausible organizations instead of one clean benchmark answer.

\subsubsection{Direct fix}\label{direct-fix-2}

\begin{enumerate}
\def\labelenumi{\arabic{enumi}.}
\tightlist
\item
  Open
  \texttt{data\_retrieval/unstructured\_retrieval/triples\_to\_rdf.py}
  and inspect the mapping for \texttt{operatedBy}.
\item
  Add an operator canonicalization table so synonyms and historical
  names resolve to the chosen ontology individual.
\item
  Restrict \texttt{operatedBy} assertions so they attach only to
  validated line entities.
\item
  Ensure structured inputs can override noisier text-derived operator
  assertions when both exist.
\item
  Re-run CQ8 and CQ14 after each change.
\end{enumerate}

\subsubsection{Files to change}\label{files-to-change-2}

\begin{itemize}
\tightlist
\item
  \texttt{data\_retrieval/unstructured\_retrieval/triples\_to\_rdf.py}
\item
  \texttt{data\_retrieval/structured\_retrieval/tfl\_api\_processor.py}
\end{itemize}

\subsection{Priority 4: Separate Night Tube from night bus evidence more
strictly}\label{priority-4-separate-night-tube-from-night-bus-evidence-more-strictly}

\subsubsection{Current issue}\label{current-issue-3}

The graph has useful night-service coverage, but the Night Tube query
pulls in buses, labels, and non-line entities because the current build
treats several night-service subjects too loosely.

\subsubsection{Direct fix}\label{direct-fix-3}

\begin{enumerate}
\def\labelenumi{\arabic{enumi}.}
\tightlist
\item
  Add a validation rule that \texttt{tfl:isNightTube\ true} may only be
  asserted on entities already typed as Underground lines.
\item
  Keep night buses in \texttt{tfl:NightBusRoute} only.
\item
  In \texttt{triples\_to\_rdf.py}, filter out generic labels such as
  \texttt{Night\ Tube}, \texttt{Night\ Bus\ network}, and route-label
  strings before type assignment.
\item
  Re-run CQ10 and compare the answer set before and after the filter.
\end{enumerate}

\subsubsection{Files to change}\label{files-to-change-3}

\begin{itemize}
\tightlist
\item
  \texttt{data\_retrieval/unstructured\_retrieval/triples\_to\_rdf.py}
\item
  optionally
  \texttt{data\_retrieval/structured\_retrieval/tfl\_api\_processor.py}
\end{itemize}

\subsection{Priority 5: Make the live extractor configurable without
editing source
code}\label{priority-5-make-the-live-extractor-configurable-without-editing-source-code}

\subsubsection{Current issue}\label{current-issue-4}

The current repository is intentionally cache-backed, which is good for
reproducibility, but local reruns are harder than they need to be
because \texttt{API\_KEY}, \texttt{model}, and \texttt{USE\_LLM} are
controlled directly in source code.

\subsubsection{Direct fix}\label{direct-fix-4}

\begin{enumerate}
\def\labelenumi{\arabic{enumi}.}
\tightlist
\item
  Read \texttt{API\_KEY}, \texttt{model}, and \texttt{USE\_LLM} from
  environment variables with sensible defaults.
\item
  Preserve the current default behaviour so the repo still runs in
  cache-backed mode out of the box.
\item
  Document the environment-variable names in the prompt documentation
  and README.
\item
  Add a short smoke-test command for local live extraction runs.
\end{enumerate}

\subsubsection{Files to change}\label{files-to-change-4}

\begin{itemize}
\tightlist
\item
  \texttt{data\_retrieval/unstructured\_retrieval/extract\_triples.py}
\item
  \texttt{docs/submission/prompt-documentation.md}
\item
  optionally the repository root \texttt{README.md}
\end{itemize}

\subsection{Priority 6: Add a post-build validation
gate}\label{priority-6-add-a-post-build-validation-gate}

\subsubsection{Current issue}\label{current-issue-5}

The project already has a competency-question harness, but it would
benefit from a second automated gate focused on graph integrity and
coursework-sensitive targets.

\subsubsection{Direct fix}\label{direct-fix-5}

\begin{enumerate}
\def\labelenumi{\arabic{enumi}.}
\tightlist
\item
  Add a validation script that checks:

  \begin{itemize}
  \tightlist
  \item
    undeclared \texttt{tfl:} classes
  \item
    presence of key benchmark instances
  \item
    counts for core classes such as operators, zones, disruptions, and
    journeys
  \item
    non-empty results for selected high-value SPARQL queries
  \end{itemize}
\item
  Run this script automatically after graph serialization in
  \texttt{main.py} or as a separate verification command.
\item
  Keep the checks simple and deterministic so they are easy to trust.
\end{enumerate}

\subsubsection{Files to change}\label{files-to-change-5}

\begin{itemize}
\tightlist
\item
  add \texttt{scripts/validate\_final\_graph.py}
\item
  \texttt{main.py}
\item
  optionally
  \texttt{competency\_questions/competency\_question\_test.py}
\end{itemize}

\subsection{Suggested Order Of Work}\label{suggested-order-of-work}

\begin{enumerate}
\def\labelenumi{\arabic{enumi}.}
\tightlist
\item
  Fix ontology alignment for emitted classes.
\item
  Tighten Night Tube and operator precision.
\item
  Add benchmark enrichment for fares, journeys, facilities, and
  disruptions.
\item
  Add the post-build validation gate.
\item
  Move live extractor configuration to environment variables.
\end{enumerate}

\subsection{Expected Payoff}\label{expected-payoff}

If these steps are implemented, the most likely benefits are:

\begin{itemize}
\tightlist
\item
  cleaner ontology-to-instance alignment
\item
  stronger competency-question precision
\item
  better coverage on missing benchmark queries
\item
  easier local reruns of the unstructured extraction stage
\item
  more confidence that future pipeline changes do not silently reduce
  graph quality
\end{itemize}

\end{document}
